\documentclass[11pt,a4paper]{article}
\usepackage[utf8]{inputenc}
\usepackage[T1]{fontenc}
\usepackage[spanish,es-noshorthands]{babel}
\usepackage[margin=2.5cm]{geometry}
\usepackage{amsmath,amssymb,mathtools}
\usepackage{enumitem}
\usepackage{booktabs}
\usepackage{tikz}
\usepackage{pgfplots}
\pgfplotsset{compat=1.18}

\newcounter{ejnum}
\newenvironment{ejercicio}{\refstepcounter{ejnum}\par\medskip\noindent\textbf{Ejercicio \theejnum.}~}{\par}
\newenvironment{solucion}{\par\noindent\textit{Solución.}~}{\par\vspace{1em}}

\title{Práctica 1: Estimación Puntual y Máxima Verosimilitud \\ \large Unidad 6: Estimación de parámetros}
\author{Probabilidad y Estadística (5765) \\ Universidad Nacional del Sur}
\date{}

\begin{document}
\maketitle

\begin{ejercicio}
Una máquina inyectora produce piezas plásticas. Se mide el peso (en gramos) de $n=6$ piezas:
\[
50.2,\ 49.8,\ 50.5,\ 49.6,\ 50.1,\ 49.9.
\]
Calcular las estimaciones puntuales de la media $\mu$ y de la varianza $\sigma^2$ usando los estimadores $\overline X$ y $S^2$ (insesgado).
\end{ejercicio}

\begin{solucion}
\[
\overline x = \frac{50.2+49.8+50.5+49.6+50.1+49.9}{6} = \frac{300.1}{6} = 50.0167 \text{ g}.
\]
Calculamos las desviaciones al cuadrado respecto de $\overline x = 50.0167$:
\[
(0.1833)^2+(-0.2167)^2+(0.4833)^2+(-0.4167)^2+(0.0833)^2+(-0.1167)^2
\]
\[
\approx 0.0336+0.0470+0.2336+0.1736+0.0069+0.0136 = 0.5083.
\]
\[
s^2 = \frac{0.5083}{n-1} = \frac{0.5083}{5} \approx 0.1017 \text{ g}^2.
\]
Estimaciones puntuales: $\widehat\mu = 50.017$ g, $\widehat{\sigma^2} = 0.1017$ g$^2$ (con $s\approx 0.319$ g).
\end{solucion}

\begin{ejercicio}
Demostrar que la varianza muestral $\widehat\sigma^2 = \frac{1}{n}\sum_{i=1}^n (X_i-\overline X)^2$ es un estimador \textbf{sesgado} de $\sigma^2$, y calcular explícitamente su sesgo para $n=10$ si $\sigma^2=4$.
\end{ejercicio}

\begin{solucion}
Sabemos (ver teoría) que $E\left[\sum_{i=1}^n(X_i-\overline X)^2\right] = (n-1)\sigma^2$. Entonces
\[
E(\widehat\sigma^2) = \frac{1}{n}(n-1)\sigma^2 = \frac{n-1}{n}\sigma^2 \neq \sigma^2 \quad (\text{salvo caso trivial}),
\]
por lo tanto $\widehat\sigma^2$ es sesgado, con sesgo
\[
B(\widehat\sigma^2) = E(\widehat\sigma^2)-\sigma^2 = \frac{n-1}{n}\sigma^2 - \sigma^2 = -\frac{\sigma^2}{n}.
\]
Para $n=10$, $\sigma^2=4$: $B(\widehat\sigma^2) = -\dfrac{4}{10} = -0.4$. En promedio, $\widehat\sigma^2$ subestima a $\sigma^2$ en $0.4$ unidades.
\end{solucion}

\begin{ejercicio}
Sea $X_1,\dots,X_n$ una m.a.\ de una población con $E(X_i)=\mu$ y $\mathrm{Var}(X_i)=\sigma^2$. Se propone el estimador
\[
\widehat\mu_1 = \frac{X_1+X_n}{2}
\]
(promedio del primer y del último dato). Analizar si $\widehat\mu_1$ es insesgado y comparar su varianza con la de $\overline X$ para $n=8$.
\end{ejercicio}

\begin{solucion}
\textbf{Insesgadez:}
\[
E(\widehat\mu_1) = \frac{E(X_1)+E(X_n)}{2} = \frac{\mu+\mu}{2} = \mu.
\]
Por lo tanto $\widehat\mu_1$ es insesgado.

\textbf{Varianza:} como $X_1$ y $X_n$ son independientes (muestra aleatoria),
\[
\mathrm{Var}(\widehat\mu_1) = \frac{1}{4}\left[\mathrm{Var}(X_1)+\mathrm{Var}(X_n)\right] = \frac{1}{4}(2\sigma^2) = \frac{\sigma^2}{2}.
\]
Para $\overline X$: $\mathrm{Var}(\overline X) = \sigma^2/n = \sigma^2/8$.

Comparando: $\mathrm{Var}(\widehat\mu_1) = \sigma^2/2 = 4\sigma^2/8 > \sigma^2/8 = \mathrm{Var}(\overline X)$.

La eficiencia relativa de $\overline X$ respecto de $\widehat\mu_1$ es $\mathrm{ER} = \dfrac{\sigma^2/2}{\sigma^2/8} = 4$: $\overline X$ es $4$ veces más eficiente. Aunque ambos son insesgados, $\overline X$ aprovecha toda la información muestral y es preferible.
\end{solucion}

\begin{ejercicio}
Calcular el error cuadrático medio (ECM) de un estimador $\widehat\theta$ tal que $E(\widehat\theta) = \theta + 2$ (sesgo constante $+2$) y $\mathrm{Var}(\widehat\theta) = 3$. Comparar con un estimador insesgado alternativo $\widehat\theta_2$ con $\mathrm{Var}(\widehat\theta_2)=6$. ¿Cuál tiene menor ECM?
\end{ejercicio}

\begin{solucion}
Sesgo de $\widehat\theta$: $B(\widehat\theta) = E(\widehat\theta)-\theta = 2$.
\[
\mathrm{ECM}(\widehat\theta) = \mathrm{Var}(\widehat\theta) + B(\widehat\theta)^2 = 3 + 2^2 = 3+4 = 7.
\]
Para $\widehat\theta_2$ (insesgado, $B=0$):
\[
\mathrm{ECM}(\widehat\theta_2) = \mathrm{Var}(\widehat\theta_2) + 0 = 6.
\]
Como $\mathrm{ECM}(\widehat\theta_2)=6 < 7=\mathrm{ECM}(\widehat\theta)$, el estimador insesgado $\widehat\theta_2$ tiene menor ECM, a pesar de tener mayor varianza: el sesgo de $\widehat\theta$ penaliza más de lo que ahorra en varianza.
\end{solucion}

\begin{ejercicio}
Probar que $\overline X$ es un estimador consistente de $\mu$ cuando $X_1,\dots,X_n$ es m.a.\ de una población con $\mathrm{Var}(X_i)=\sigma^2 < \infty$. Luego, usando la desigualdad de Chebyshev, acotar $P(|\overline X-\mu|>0.5)$ si $\sigma^2=9$ y $n=100$.
\end{ejercicio}

\begin{solucion}
$\overline X$ es insesgado ($E(\overline X)=\mu$) y $\mathrm{Var}(\overline X) = \sigma^2/n \to 0$ cuando $n\to\infty$. Por lo tanto $\mathrm{ECM}(\overline X)=\mathrm{Var}(\overline X)\to 0$, lo que implica consistencia (convergencia en probabilidad a $\mu$).

Aplicando Chebyshev con $\varepsilon=0.5$:
\[
P(|\overline X-\mu|>0.5) \le \frac{\mathrm{Var}(\overline X)}{0.5^2} = \frac{\sigma^2/n}{0.25} = \frac{9/100}{0.25} = \frac{0.09}{0.25} = 0.36.
\]
Es decir, la probabilidad de que $\overline X$ se aleje de $\mu$ en más de $0.5$ unidades es, a lo sumo, $0.36$.
\end{solucion}

\begin{ejercicio}
Un control de calidad registra si cada artículo producido es defectuoso ($X_i=1$) o no ($X_i=0$), modelado como Bernoulli($p$). En una muestra de $n=150$ artículos se detectan $18$ defectuosos. Hallar el estimador de máxima verosimilitud de $p$ y de la probabilidad de \textbf{no} tener un defectuoso, $q=1-p$, usando la propiedad de invarianza.
\end{ejercicio}

\begin{solucion}
Para Bernoulli, el EMV es $\widehat p = \overline X = $ proporción de éxitos (defectuosos):
\[
\widehat p = \frac{18}{150} = 0.12.
\]
Por invarianza del EMV, como $g(p)=1-p$ es biunívoca,
\[
\widehat q = 1-\widehat p = 1-0.12 = 0.88.
\]
Estimación: el $12\%$ de los artículos son defectuosos, y el $88\%$ no lo son.
\end{solucion}

\begin{ejercicio}
El tiempo (en horas) hasta la falla de un componente electrónico se modela como $X\sim\text{Exp}(\lambda)$. Se registran los tiempos de falla de $5$ componentes:
\[
120,\ 95,\ 150,\ 80,\ 105.
\]
Obtener la estimación de máxima verosimilitud de $\lambda$ y del tiempo medio de vida $E(X)=1/\lambda$.
\end{ejercicio}

\begin{solucion}
Sabemos que el EMV de $\lambda$ en la exponencial es $\widehat\lambda = 1/\overline x$.
\[
\overline x = \frac{120+95+150+80+105}{5} = \frac{550}{5} = 110 \text{ horas}.
\]
\[
\widehat\lambda = \frac{1}{110} \approx 0.00909 \text{ fallas/hora}.
\]
Por invarianza, el EMV de $E(X)=1/\lambda$ es
\[
\widehat{E(X)} = \frac{1}{\widehat\lambda} = \overline x = 110 \text{ horas}.
\]
\end{solucion}

\begin{ejercicio}
Se midieron las estaturas (en cm) de $n=7$ estudiantes, que se suponen $N(\mu,\sigma^2)$:
\[
168,\ 172,\ 165,\ 178,\ 170,\ 174,\ 169.
\]
Hallar las estimaciones de máxima verosimilitud de $\mu$ y $\sigma^2$.
\end{ejercicio}

\begin{solucion}
El EMV de $\mu$ es $\widehat\mu=\overline x$:
\[
\overline x = \frac{168+172+165+178+170+174+169}{7} = \frac{1196}{7} = 170.857 \text{ cm}.
\]
El EMV de $\sigma^2$ es $\widehat\sigma^2 = \frac{1}{n}\sum(x_i-\overline x)^2$ (con $n$, no $n-1$). Calculamos las desviaciones respecto de $170.857$:
\begin{align*}
(168-170.857)^2 &\approx 8.163\\
(172-170.857)^2 &\approx 1.306\\
(165-170.857)^2 &\approx 34.306\\
(178-170.857)^2 &\approx 51.020\\
(170-170.857)^2 &\approx 0.735\\
(174-170.857)^2 &\approx 9.878\\
(169-170.857)^2 &\approx 3.449
\end{align*}
Suma $\approx 108.857$. Entonces
\[
\widehat\sigma^2 = \frac{108.857}{7} \approx 15.551 \text{ cm}^2, \qquad \widehat\sigma \approx 3.944 \text{ cm}.
\]
\end{solucion}

\begin{ejercicio}
El número de accidentes semanales en una planta industrial sigue una distribución $\text{Poisson}(\lambda)$. Se observaron los siguientes registros en $10$ semanas:
\[
2,\ 0,\ 1,\ 3,\ 2,\ 1,\ 0,\ 2,\ 4,\ 1.
\]
a) Hallar el EMV de $\lambda$. b) Usando invarianza, estimar $P(X=0)=e^{-\lambda}$, la probabilidad de una semana sin accidentes.
\end{ejercicio}

\begin{solucion}
a) El EMV de $\lambda$ en Poisson es $\widehat\lambda=\overline x$.
\[
\sum x_i = 2+0+1+3+2+1+0+2+4+1 = 16, \qquad \widehat\lambda = \frac{16}{10} = 1.6.
\]
b) Por invarianza,
\[
\widehat{P(X=0)} = e^{-\widehat\lambda} = e^{-1.6} \approx 0.2019.
\]
Se estima que, aproximadamente, el $20.19\%$ de las semanas no tendrán accidentes.
\end{solucion}

\end{document}
